% ═══════════════════════════════════════════════════════════════════════════
% FROZEN OREGON PLEADING LAYOUT — DO NOT "IMPROVE" WITHOUT HUMAN APPROVAL
% ═══════════════════════════════════════════════════════════════════════════
%
% Version:     1.0 (locked 2026-07-13)
% Provenance:  Matches the LaTeX look of the golden Motion PDF dated
%              2026-07-10 (pdfTeX 1.40.28 / TeX Live 2025), archived at:
%              templates/golden/Motion_to_Set_Aside_Void_Judgment_GOLDEN_2026-07-10.pdf
%
% Source ancestor (pre-mutation 32-line municipal draft):
%   templates/golden/municipal_oregon_city_consolidated_LEGACY_SOURCE.tex
%   (was living at ~/Downloads/municipal_oregon_city_consolidated.tex — fragile)
%
% Locked features (the ones Annika fought to get right):
%   • 28-line pleading grid (not 32)
%   • Line-number baseline shift +3.5pt so 1–28 track body text
%   • Double vertical margin rule at 1.3in and 1.265in
%   • Letter page, 1in top/bottom, 1.5in left, 1in right
%   • Caption shell (injected by compile_pleading): 10.5pt single-spaced
%     parties left / case no + title right, tight 0.25in gap, rule touching center
%
% Antiregression:
%   • compile_pleading.py MUST prefer this file over Downloads/
%   • ./do layout-status verifies this path + golden PDF exist
%   • Do not reintroduce runtime mutations that rewrite 32→28 unless
%     loading the LEGACY_SOURCE only.
%
% Footer title string is a PLACEHOLDER replaced at compile time.
% ═══════════════════════════════════════════════════════════════════════════
\documentclass[12pt]{article}
\usepackage[letterpaper, top=1in, bottom=1in, left=1.5in, right=1in,
            headheight=0pt, headsep=0pt, footskip=0.4in]{geometry}
\usepackage{lmodern}
\usepackage{ragged2e}
\usepackage{calc}
\usepackage{eso-pic}
\usepackage{fancyhdr}

% Keep headings with following text (replaces needspace.sty — not always installed)
\makeatletter
\newcommand{\Needspace}[1]{%
  \par \penalty-100 \begingroup
  \setlength{\dimen@}{#1}%
  \dimen@ii\pagegoal \advance\dimen@ii-\pagetotal
  \ifdim \dimen@ii<\dimen@
    \ifdim \dimen@ii>0pt
      \newpage
    \fi
  \fi
  \endgroup
}
\makeatother

% =====================================================
% 28-LINE PLEADING GRID — FROZEN GOLDEN LAYOUT
% =====================================================
\newlength{\gridline}
\setlength{\gridline}{\textheight/28}   % 9in / 28 — locked to golden Motion PDF 2026-07-10
\setlength{\topskip}{\gridline}

\makeatletter
\newcommand{\uld}[1]{\strip@pt\dimexpr#1\relax}
\makeatother

\newcount\gridcnt
\AddToShipoutPictureBG{%
  \AtPageUpperLeft{%
    \setlength{\unitlength}{1pt}%
    \gridcnt=0
    \loop
      \advance\gridcnt by 1
      \put(\uld{1.05in},\uld{-1in-\gridline*\gridcnt+3.5pt}){%
        \makebox(0,0)[r]{\footnotesize\the\gridcnt}}%
    \ifnum\gridcnt<28
    \repeat
    \put(\uld{1.3in},\uld{-1in}){\line(0,-1){\uld{\gridline*28}}}\put(\uld{1.265in},\uld{-1in}){\line(0,-1){\uld{\gridline*28}}}%
  }%
}

\pagestyle{fancy}
\fancyhf{}
\renewcommand{\headrulewidth}{0pt}
\renewcommand{\footrulewidth}{0pt}
\fancyfoot[L]{\footnotesize Defendant's Consolidated Petition for Hardship Relief, Page \thepage}

% =====================================================
% Typography & paragraph control
% =====================================================
\setlength{\parskip}{0pt}
\setlength{\parindent}{0.5in}
\setlength{\topsep}{0pt}
\setlength{\partopsep}{0pt}
\raggedbottom
\raggedright              % prevent overfull hboxes, cleaner pro se look
\widowpenalty=10000        % disallow widows
\clubpenalty=10000         % disallow orphans
\brokenpenalty=10000       % discourage breaks after hyphens
\hyphenpenalty=300         % reduce hyphenation frequency
\exhyphenpenalty=300
\emergencystretch=1em      % allow slight line stretch to avoid overfulls

\begin{document}
% Body is fully generated by compile_pleading.py (caption shell + content).
% This file exists so the PREAMBLE is never lost or re-derived from Downloads.
\romanhead{IN THE COURT OF APPEALS OF THE STATE OF OREGON}

\para{1}{ANNIKA ERIKSSON, Appellant, v. THE BANK OF NEW YORK MELLON FKA THE BANK OF NEW YORK, AS TRUSTEE FOR THE CERTIFICATEHOLDERS OF CWALT, INC., ALTERNATIVE LOAN TRUST 2006-OC10, MORTGAGE PASS-THROUGH CERTIFICATES, SERIES 2006-OC10, and THE BANK OF NEW YORK MELLON FKA THE BANK OF NEW YORK, AS TRUSTEE FOR THE CERTIFICATEHOLDERS OF CWALT, INC., ALTERNATIVE LOAN TRUST 2006-OC11, MORTGAGE PASS-THROUGH CERTIFICATES, SERIES 2006-OC11; NEWREZ, LLC D/B/A SHELLPOINT MORTGAGE SERVICING; and ALL OTHER UNKNOWN PARTIES, Respondents.}

\para{2}{Appellate Case No. A191649 | Clackamas County Circuit Court Case No. 24CV21417}

\subhead{NOTICE OF SUPPLEMENTAL STATUS UPDATE REGARDING MOTION TO SET ASIDE JUDGMENT IN CASE NO. 17CV38591, AND MOTION FOR INTERIM RELIEF PENDING RULING ON EMERGENCY MOTION FOR STAY}

\para{3}{Appellant Annika Eriksson, Pro Se, respectfully submits this Notice to inform the Court of a material change in status occurring after the filing of Appellant's Emergency Motion Under ORAP 7.35 for Stay of Enforcement Pending Appeal, filed August 10, 2026. Appellant further moves, as set forth below, for interim relief preserving the status quo while the Emergency Motion remains under submission.}

\para{4}{In her Declaration in support of the Emergency Motion, Appellant stated that her Motion to Set Aside Void Judgment in Clackamas County Circuit Court Case No. 17CV38591 (Envelope \#14448055) was submitted on August 5, 2026, and was, at that time, under review by the trial court clerk.}

\para{5}{Appellant now advises the Court that, as of the date of this filing, the Motion to Set Aside Judgment (Envelope \#14448055) has been Accepted by the Clackamas County Circuit Court and is now formally entered on the docket in Case No. 17CV38591.}

\para{6}{This development is material to the pending Emergency Motion because the Motion to Set Aside Judgment directly challenges the validity of the 2019 foreclosure judgment that underlies the reformation judgment presently on appeal in this matter. The Motion's acceptance onto the trial court docket confirms that a live, pending challenge to the 2019 judgment now exists in the trial court.}

\para{7}{Appellant respectfully requests that the Court take note of this development in its consideration of the pending Emergency Motion for Stay of Enforcement Pending Appeal.}

\para{8}{Appellant further advises the Court that, as of the date of this filing, no interim or temporary order addressing the Emergency Motion has yet been entered. Respondents' response is not due until August 24, 2026, and the General Judgment on appeal expressly authorizes Respondents to record a reformed sheriff's deed at any time. Because recording could occur at any moment during the pendency of the response period, and because a live, docketed challenge to the validity of the underlying 2019 judgment now exists in Case No. 17CV38591, Appellant respectfully requests that the Court consider whether interim relief preserving the status quo is warranted while the Emergency Motion remains under submission.}

\para{9}{The standard governing the Emergency Motion is set out in ORS 19.350(3)(a), under which the Court considers: (a) the likelihood that Appellant will prevail on appeal; (b) whether the appeal is taken in good faith and not for the purpose of delay; (c) whether there is support for the appeal in law or in fact; and (d) the nature of harm to Appellant, to other parties, and to the public that will likely result from the grant or denial of a stay. See Neal v. Neal, Court of Appeals No. A169261 (Order Granting Stay, Jan. 17, 2020) (unpublished order of the Appellate Commissioner).}

\para{10}{In Neal, arising from Clackamas County Circuit Court No. 18CV02117, the Appellate Commissioner granted a stay under materially analogous circumstances. The appellant there had been dismissed from the trial court on procedural grounds rather than on the merits of her claims to real property, which included quiet title. The Commissioner held that a preliminary showing of a reasonable likelihood of prevailing was sufficient to satisfy the first factor, and found that a party's extensive prior litigation history did not, standing alone, establish bad faith or support denial of a stay. The Commissioner further held that denial of a stay would leave the opposing party "free to dispose of the property," causing irreparable harm should the appellant ultimately prevail, and accordingly stayed the judgment specifically to prevent conveyance, transfer, sale, or encumbrance of the property pending appeal.}

\para{11}{The same considerations apply here with equal, if not greater, force. Appellant's now-docketed Motion to Set Aside Void Judgment in Case No. 17CV38591 directly challenges the validity of the 2019 foreclosure judgment underlying the reformation judgment on appeal. As in Neal, a live and pending challenge to the foundational judgment is itself a preliminary showing sufficient to support the first factor, particularly because a ruling in Appellant's favor in Case No. 17CV38591 would necessarily undermine the reformation judgment presently before this Court. Nor does the existence of related, interconnected filings across Case Nos. 17CV38591 and 24CV21417 support an inference of delay or bad faith; as in Neal, the interrelated nature of these proceedings reflects the interconnected legal defects at issue, not tactical delay.}

\para{12}{The irreparable-harm factor is directly controlled by Neal. There, as here, the harm at issue was not compensable in damages: it was the risk that the opposing party could alter the status of real property before the appeal, and any related challenge to the underlying judgment, could be resolved. Here, the General Judgment on appeal expressly authorizes Respondents to record a reformed sheriff's deed at any time, with no restriction pending either the response period or this Court's ruling on the Emergency Motion. If the reformed deed is recorded before this Court rules, Appellant's appeal, and her pending motion in Case No. 17CV38591, could be rendered effectively moot as to third-party title questions — precisely the irreparable-harm scenario that warranted a stay in Neal.}

\para{13}{Given the imminent risk that recording may occur before Respondents' response is due or before the Court has ruled on the Emergency Motion, Appellant respectfully requests that the Court act on an interim basis, in advance of full briefing, to preserve the status quo.}

\romanhead{MOTION FOR INTERIM RELIEF}

\para{14}{For the foregoing reasons, Appellant respectfully moves this Court for a temporary order enjoining Respondents from recording the reformed sheriff's deed, or otherwise conveying, transferring, encumbering, or altering title to the subject property, pending this Court's ruling on the Emergency Motion for Stay of Enforcement Pending Appeal, or further order of this Court.}

\para{15}{DATED: August 17, 2026.}

\para{16}{Respectfully submitted,}

\para{17}{_________________________________}

\para{18}{Annika Eriksson, Appellant Pro Se}

\para{19}{P.O. Box 1108 (ACP 0030-24)}

\para{20}{Salem, OR 97308}

\para{21}{(971) 359-8578}

\para{22}{eriksson@icloud.com}

\romanhead{CERTIFICATE OF SERVICE}

\para{23}{I certify that on August 17, 2026, I served a true and correct copy of the foregoing Notice of Supplemental Status Update and Motion for Interim Relief on counsel of record by First-Class Mail and e-mail courtesy copy: Gregor A. Hensrude \& Chancellor K. Eagle, Klinedinst PC (ghensrude@klinedinstlaw.com; ceagle@klinedinstlaw.com) by First-Class Mail \& e-mail courtesy copy.}

\para{24}{DATED: August 17, 2026.}

\para{25}{_________________________________}

\para{26}{Annika Eriksson, Appellant Pro Se}
\end{document}
